
% ============================================================
% METODOLOGÍA DE APRENDIZAJE
%
% Esta sección combina:
% - párrafos normales;
% - una lista descriptiva;
% - una lista secundaria;
% - una nota final.
% ============================================================

\section{Metodología de aprendizaje}

El curso se articula principalmente en torno a la modalidad de
\textit{Charlas Magistrales Semi-Dirigidas}. Bajo este formato, las
sesiones no consisten en la exposición unilateral de contenidos
por parte del docente, sino en dinámicas de interacción orientadas
donde el docente actúa como guía que contextualiza los problemas nodales,
modela la precisión terminológica y conduce la discusión crítica
entre y con los estudiantes.

Para llevar a cabo esta propuesta, cada sesión de clase se organiza mediante 
dos momentos complementarios:

\begin{description}[
    style=nextline,
    leftmargin=0pt,
    labelindent=0pt,
    labelwidth=0pt,
    labelsep=0pt,
    itemsep=7pt
]

\item[\textbf{Charlas Magistrales Semi-Dirigidas}]
Constituyen el primer bloque de cada sesión. En este espacio, con una 
participación e interacción mayoritaria del docente, se presentan, 
contextualizan y guían los conceptos lingüísticos fundamentales del tema programado, 
manteniendo un diálogo abierto con los estudiantes.  

\item[\textbf{Espacio de Aplicación y Reflexión}]
Corresponde al segmento final de la sesión, desarrollado tras el bloque magistral. 
Se orienta a la resolución de casos concretos de análisis de datos o al debate de preguntas 
guiadas sobre lo expuesto. La participación activa, continua y bien fundamentada en estos espacios
constituirá el insumo central para la evaluación del componente de participación.

\end{description}

En complemento a la dinámica habitual de las sesiones, el trabajo práctico, analítico y evaluativo del 
curso se articula a través de los siguientes componentes:

\begin{description}[
    style=nextline,
    leftmargin=0pt,
    labelindent=0pt,
    labelwidth=0pt,
    labelsep=0pt,
    itemsep=7pt
]

\item[\textbf{Talleres Conceptuales Grupales (4)}]
Evaluaciones de aplicación práctica distribuidas estratégicamente a lo largo del semestre. 
Cada taller agrupa los contenidos abordados en bloques de aproximadamente tres sesiones para que los 
grupos consoliden los conceptos vistos en clase mediante ejercicios de análisis, argumentación y/o reflexión, 
usualmente con casos reales del español, la LSC u otras lenguas de señas cuando el ejercicio así lo requiera.

\item[\textbf{Trabajo Final Integrador: Recurso Pedagógico Bimodal}]
Elaboración grupal de un recurso pedagógico, visual y esquemático que presente una mirada integrada y comparativa
de los principales contenidos fonológicos y morfológicos abordados durante el curso en español y Lengua de Señas 
Colombiana (LSC). El recurso deberá:

\begin{itemize}[leftmargin=*,itemsep=3pt,topsep=4pt]
    \item \textbf{Organizar} los contenidos a partir de cuatro acciones analíticas: \textbf{observar}, \textbf{comparar}, \textbf{sustentar} y \textbf{explicar}.
    \item Establecer relaciones, \textbf{semejanzas y diferencias estructurales} entre ambas modalidades lingüísticas.
    \item \textbf{Profundizar en un problema lingüístico} específico seleccionado por el grupo.
    \item Estar dirigido estratégicamente a estudiantes e interesados en el campo de la traducción e interpretación, 
    incorporando una reflexión explícita sobre la manera en que la comprensión lingüística de ambas lenguas aporta a 
    la formación y futura práctica del traductor e intérprete.
\end{itemize}

Se privilegiarán formas de presentación pedagógicas, accesibles y altamente visuales que permitan comunicar conceptos 
lingüísticos con claridad y rigor analítico.

\end{description}

\textbf{Nota.}
La elaboración de los talleres grupales y el desarrollo del recurso pedagógico final se realizarán en grupos de trabajo 
de máximo cuatro integrantes, promoviendo el aprendizaje colaborativo en el análisis de datos.
